\documentclass{ximera}
\title{Worksheet, Simple Harmonic Motion}


\newcommand{\pskip}{\vskip 0.1 in}

\begin{document}
\begin{abstract}
Simple harmonic motion worksheet.
\end{abstract}
\maketitle

\emph{Directions} Use the protractor below to help find approximate answers to the questions that follow. Note the unusual orientation of the coordinate axes. Do \emph{not} use the sine or cosine functions.


\begin{onlineOnly}
    \begin{center}
\desmos{9xuhmux5qr}{900}{600}
\end{center}
\end{onlineOnly}

\href{https://www.desmos.com/calculator/9xuhmux5qr}{142: Revolution Protractor 2}


\begin{question} \label{Qpdf30r3mmm}

A mass hanging from a spring oscillates in simple harmonic motion about its equilibrium position between displacements $s=-16$cm and $s=16$cm from equilibrium with a period of $4$ seconds. The mass is released from rest at $s=-16$cm.

A point $P$ on the circle of radius $16cm$ centered at the equilibrium position drives the simple harmonic motion. The point $Q$ on the circle of radius $5$cm moves in tandem with $P$, having the same polar angle (meausred in revolutions from the positive $x$-axis) at all times. 

\begin{onlineOnly}
    \begin{center}
\desmos{tskmnyezxz}{900}{600}
\end{center}
\end{onlineOnly}

\href{https://www.desmos.com/calculator/tskmnyezxz}{142: Vertical Spring}

\begin{enumerate}
\item Find a function 
\[
    \theta = a(t) \, , \, t \geq 0 ,
\]
that expresses the polar angles of $P$ and $Q$, measured counterclockwise from the positive $x$-axis in \emph{revolutions}, in terms of the number of seconds since the mass was released.

\item Find a function
\[
    s = f(x) \, , \, -5 \leq x \leq 5 ,
\]
that expresses the displacement from equilibrium of the oscillating mass in terms of the $x$-coordinate of point $Q$.

\item Let $x = g(\theta)$, $\theta \geq 0$, be the function that expresses the $x$-coordinate of point $Q$ in terms of the polar angle $\theta$. Express the function
\[
   s = h(t) \, , \, t\geq 0,
\]
that gives the position of the mass relative to its equilibrium position (in cm) in terms of the number of seconds since the mass was released in terms of the functions $a$, $f$, and $g$. Work in general. Do \emph{not} use your specific expressions for the functions $a$ and $f$.

\item Sketch a rough graph of the function $s=f(t)$, $t\geq 0$. Label the coordinates of at least four key points on the graph.

\item Use the protractor and your functions from parts (a) and (b) to find approximate answers to each of the following questions.

\begin{enumerate}
\item Find the position of the mass at time $t=1.75$ seconds after being released.

\item Approximate the first two times when the mass is at position $s=-9$cm.

\item Approximate the twentieth time the mass is at position $s=3$cm.
\end{enumerate}

\end{enumerate}

\end{question}



\begin{question} \label{Qdlferermm}

Between consecutive midnights the depth of the water at Edmond's pier oscillates in simple harmonic motion about its mean, reaching a maximum depth of $17$ feet at 2:00am and a minimum depth of $8$ feet at 7:30am. 
 
\begin{enumerate}
\item Repeat parts (a)-(d) of Question 1, appropriately modified.

\item Approximate the water's depth at 5am and at 3pm.

\item Approximate all times when the water is $15$ feet deep.

\item Approximate the depth of the water $30$ minutes after the depth is $10$ feet. Give all possibilities.
\end{enumerate}
\end{question}




\end{document}